%==============================================================================
%  A Rigorous Foundation of Gauss--Hermite and Gauss--Legendre Quadrature
%  with Applications to the Generalized Interpolation Lattice Boltzmann Method
%
%  Author : P.-C.
%  Date   : 27 June 2026
%  Version: 1.0
%
%  Compile with:  pdflatex -> bibtex -> pdflatex x 2
%  (no external .bib needed; uses thebibliography environment)
%==============================================================================
\documentclass[11pt,a4paper]{article}

%---------- Page geometry ----------
\usepackage[margin=1in]{geometry}

%---------- Encoding & fonts ----------
\usepackage[T1]{fontenc}
\usepackage[utf8]{inputenc}
\usepackage{lmodern}

%---------- Math packages ----------
\usepackage{amsmath,amssymb,amsthm,amsfonts}
\usepackage{mathtools}
\usepackage{bm}

%---------- Tables & graphics ----------
\usepackage{booktabs}
\usepackage{array}
\usepackage{graphicx}
\usepackage{multirow}

%---------- Colour ----------
\usepackage[dvipsnames]{xcolor}

%---------- Hyperref (load last among most) ----------
\usepackage[
  colorlinks=true,
  linkcolor=blue!60!black,
  citecolor=blue!60!black,
  urlcolor=blue!70!black,
  pdftitle={A Rigorous Foundation of Gauss--Hermite and Gauss--Legendre Quadrature
            with Applications to the Generalized Interpolation Lattice Boltzmann Method},
  pdfauthor={P.-C.}
]{hyperref}

%---------- Theorem environments ----------
\newtheorem{theorem}{Theorem}[section]
\newtheorem{lemma}[theorem]{Lemma}
\newtheorem{proposition}[theorem]{Proposition}
\newtheorem{corollary}[theorem]{Corollary}

\theoremstyle{definition}
\newtheorem{definition}[theorem]{Definition}
\newtheorem{remark}[theorem]{Remark}
\newtheorem{observation}[theorem]{Observation}

%---------- Custom math macros ----------
\newcommand{\R}{\mathbb{R}}
\newcommand{\N}{\mathbb{N}}
\newcommand{\PP}{\mathbb{P}}
\newcommand{\inner}[1]{\left\langle #1 \right\rangle}
\newcommand{\ip}[2]{\langle #1,\, #2 \rangle}
\newcommand{\ipw}[2]{\langle #1,\, #2 \rangle_{w}}
\newcommand{\I}{\mathcal{I}}
\newcommand{\Q}{\mathcal{Q}}
\newcommand{\HE}{\mathrm{He}}
\newcommand{\GL}{\mathrm{GL}}
\newcommand{\GH}{\mathrm{GH}}
\newcommand{\eq}{\mathrm{eq}}
\newcommand{\MB}{\mathrm{MB}}
\DeclareMathOperator{\supp}{supp}
\DeclareMathOperator{\diag}{diag}
\DeclareMathOperator{\sgn}{sgn}

%---------- Section heading shaping ----------
\usepackage{titlesec}
\titleformat{\section}{\large\bfseries}{\thesection.}{0.6em}{}
\titleformat{\subsection}{\normalsize\bfseries}{\thesubsection}{0.5em}{}

%==============================================================================
\title{\Large\bfseries
A Rigorous Foundation of Gauss--Hermite and Gauss--Legendre Quadrature
with Applications to the Generalized Interpolation Lattice Boltzmann Method}

\author{P.-C.\thanks{Correspondence: \texttt{s8313697@gmail.com}}}
\date{27 June 2026}

%==============================================================================
\begin{document}

\maketitle

\begin{abstract}
\noindent
We present a unified, rigorous theoretical framework for the Gauss--Hermite and
Gauss--Legendre quadrature rules and elucidate their complementary roles within
the Generalized Interpolation Lattice Boltzmann Method (GILBM). Working in
weighted $L^2$ spaces, we recover the orthogonality and completeness structure
of the classical orthogonal polynomial families
\cite{Szego1975} and provide a self-contained five-step proof of the Gaussian
quadrature precision theorem, establishing algebraic exactness of degree
$2n-1$ for any $n$-point rule whose nodes coincide with the roots of an
orthogonal polynomial of degree $n$
\cite{Stroud1971,DavisRabinowitz1984}. By carefully distinguishing three
structural roles---the \emph{target functional} $\ipw{f}{1}$, the
\emph{quadrature nodes} (polynomial roots), and the
\emph{orthogonality-induced annihilation mechanism}---we correct a common
conceptual conflation in the applied literature.

In the kinetic-theoretic setting we prove that the Hermite weight function is
the \emph{unique} member of the classical orthogonal families that coincides
with the dimensionless Maxwell--Boltzmann distribution; this isomorphism is the
structural reason Gauss--Hermite quadrature constitutes the canonical
discretization of velocity space in the lattice Boltzmann method
\cite{HeLuo1997a,HeLuo1997b,ShanHe1998,ShanYuanChen2006}. Conversely, the
physical-space reconstruction step intrinsic to GILBM corresponds to an
unweighted integration on the unit cell and is therefore naturally governed by
Gauss--Legendre quadrature. We further show, as a consequence of the
completeness-interval theorem, that the two quadrature schemes are
\emph{not interchangeable}: each is rigidly tied to the support of its own
weight measure. The framework provides a rigorous foundation for the design of
discrete velocity sets, interpolation stencils, and a~priori error estimates in
higher-order GILBM solvers.

\medskip
\noindent\textbf{Keywords:} Gauss--Hermite quadrature; Gauss--Legendre
quadrature; orthogonal polynomials; weighted $L^2$ spaces; lattice Boltzmann
method; generalized interpolation; Maxwell--Boltzmann distribution.

\medskip
\noindent\textbf{MSC 2020:} 65D32, 33C45, 76P05, 76M28.
\end{abstract}

%==============================================================================
\section{Introduction}\label{sec:intro}

The lattice Boltzmann method (LBM) has, since the seminal work of He and Luo
\cite{HeLuo1997a}, been elevated from a heuristic discrete-kinetic model into a
rigorous Galerkin-type discretization of the continuous Boltzmann equation. The
mathematical machinery underpinning this elevation rests on two pillars:
(i)~the Hermite-polynomial expansion of the equilibrium distribution about the
local Maxwellian, and (ii)~the Gauss--Hermite quadrature that exactly preserves
a finite hierarchy of velocity moments
\cite{ShanHe1998,Philippi2006,ShanYuanChen2006}.

In the Generalized Interpolation Lattice Boltzmann Method (GILBM), the
requirement to reconstruct the post-collision distribution at off-lattice
locations introduces a \emph{second} quadrature problem, defined this time on a
bounded local cell with unit weight. This second problem is naturally
addressed by Gauss--Legendre quadrature
\cite{HeLuoDembo1996,Kruger2017,Succi2018}. The two quadratures, far from being
interchangeable, occupy distinct mathematical niches dictated by the
completeness intervals of their underlying orthogonal-polynomial families.

The present paper has three objectives:
\begin{enumerate}
  \item[(i)] \textbf{Mathematical:} to provide a self-contained, rigorous
        derivation of Gauss--Hermite and Gauss--Legendre quadrature within a
        unified weighted-$L^2$ framework, including the algebraic-precision
        theorem with a complete proof and an explicit optimality argument.
  \item[(ii)] \textbf{Conceptual:} to dispel a common misconception that the
        Gaussian quadrature integral represents an inner product against the
        orthogonal polynomial $\phi_n$. We show that the \emph{target}
        functional is $\ipw{f}{1}$ and that $\phi_n$ acts only through the
        \emph{node placement} and the \emph{orthogonality annihilation}
        mechanism.
  \item[(iii)] \textbf{Applied:} to formulate the canonical assignment of
        Gauss--Hermite to velocity space and Gauss--Legendre to physical space
        within GILBM, and to prove (via the completeness-interval theorem) that
        this assignment is unique up to trivial rescaling.
\end{enumerate}

%==============================================================================
\section{Mathematical Preliminaries}\label{sec:prelim}

\subsection{Weighted $L^2$ spaces}

Let $I\subseteq\R$ be an open interval (possibly unbounded) and
$w\colon I\to\R_{\ge 0}$ a measurable weight function satisfying
\begin{description}\itemsep0pt
  \item[(W1)] $w(x)\ge 0$ almost everywhere on $I$;
  \item[(W2)] $\int_K w(x)\,dx<\infty$ for every compact $K\subset I$;
  \item[(W3)] all moments are finite:
              $\mu_k := \int_I x^k\,w(x)\,dx<\infty$ for all $k\in\N_0$.
\end{description}

The weighted Hilbert space
\[
  L^2(I,w) \;=\;
  \left\{\, f\colon I\to\R \;\middle|\; \int_I |f(x)|^2\,w(x)\,dx<\infty\,\right\}
\]
is endowed with the inner product and norm
\[
  \ipw{f}{g} \;:=\; \int_I f(x)\,g(x)\,w(x)\,dx,
  \qquad
  \|f\|_w \;:=\; \sqrt{\ipw{f}{f}}.
\]

\subsection{Classical orthogonal polynomials}

Applying Gram--Schmidt orthogonalization to $\{1,x,x^2,\ldots\}$ inside
$L^2(I,w)$ produces a family $\{\phi_n\}_{n=0}^\infty$ of polynomials with
$\deg\phi_n=n$ and
\[
  \ipw{\phi_n}{\phi_m} \;=\; h_n\,\delta_{nm},\qquad h_n>0.
\]
The four classical families relevant here are summarized in
Table~\ref{tab:classical}.

\begin{table}[ht]
\centering
\caption{Classical orthogonal polynomial families.}
\label{tab:classical}
\begin{tabular}{lccc}
\toprule
Family & $I$ & $w(x)$ & $h_n$ \\
\midrule
Legendre  $P_n$            & $[-1,1]$        & $1$
  & $\dfrac{2}{2n+1}$ \\[4pt]
Chebyshev (1st) $T_n$      & $[-1,1]$        & $(1-x^2)^{-1/2}$
  & $\pi/2\;(n\ge1),\;\pi\;(n=0)$ \\[4pt]
Laguerre $L_n$             & $[0,\infty)$    & $e^{-x}$
  & $1$ \\[4pt]
Hermite (probabilist) $\HE_n$ & $\R$         & $(2\pi)^{-1/2}\,e^{-x^2/2}$
  & $n!$ \\[4pt]
Hermite (physicist) $H_n$  & $\R$            & $e^{-x^2}$
  & $\sqrt{\pi}\,2^n n!$ \\
\bottomrule
\end{tabular}
\end{table}

\begin{theorem}[Completeness; \cite{Szego1975}, Thm.~5.7.1]\label{thm:complete}
If $(I,w)$ satisfies \emph{(W1)--(W3)} together with the determinacy of the
Hamburger moment problem, then $\{\phi_n\}_{n=0}^\infty$ is dense in
$L^2(I,w)$, i.e.\ it constitutes a complete orthogonal basis.
\end{theorem}

Both Legendre and Hermite families satisfy the determinacy hypothesis. For
Hermite this follows from Carleman's criterion since
$\sum_n \mu_{2n}^{-1/(2n)}=+\infty$.

\begin{remark}[Completeness interval]\label{rmk:completeness}
The interval $I$ on which completeness is asserted is \emph{intrinsic} to the
polynomial family. Any quadrature derived from $\{\phi_n\}$ inherits this
interval as its natural domain; departing from $I$ destroys both orthogonality
and the precision guarantees built upon it.
\end{remark}

\subsection{Legendre and Hermite polynomials}

\paragraph{Legendre.} Three-term recurrence:
\[
  P_0(x)=1,\quad P_1(x)=x,\qquad
  (n+1)P_{n+1}(x)=(2n+1)\,x\,P_n(x)-n\,P_{n-1}(x).
\]
Rodrigues' formula:
\(
  P_n(x)=\frac{1}{2^n n!}\frac{d^n}{dx^n}\bigl[(x^2-1)^n\bigr].
\)
Generating function \cite{WhittakerWatson1927}:
\[
  G_P(x,t)\;=\;\frac{1}{\sqrt{1-2xt+t^2}}\;=\;\sum_{n=0}^\infty P_n(x)\,t^n,
  \qquad |t|<1.
\]

\paragraph{Hermite (probabilist).} Three-term recurrence:
\[
  \HE_0(x)=1,\quad \HE_1(x)=x,\qquad
  \HE_{n+1}(x)=x\,\HE_n(x)-n\,\HE_{n-1}(x).
\]
Rodrigues' formula:
\(
  \HE_n(x)=(-1)^n\,e^{x^2/2}\,\frac{d^n}{dx^n}\,e^{-x^2/2}.
\)
Generating function:
\[
  G_{\HE}(x,t)\;=\;\exp\!\bigl(xt-\tfrac{1}{2}t^2\bigr)
  \;=\;\sum_{n=0}^\infty \HE_n(x)\,\frac{t^n}{n!}.
\]

\begin{observation}[Singularity of Hermite]\label{obs:hermite-singular}
Among all classical orthogonal polynomial families, the Hermite family is the
\emph{unique} one whose orthogonality weight
$w(x)=(2\pi)^{-1/2}e^{-x^2/2}$ is itself the standard Gaussian
(Gauss-core) density. This isomorphism is the structural origin of Hermite's
irreplaceability in kinetic-theoretic discretization.
\end{observation}

\subsection{Roots}

\begin{theorem}[Real, simple, interior roots; \cite{Szego1975}, Thm.~3.3.1]
\label{thm:roots}
If $\phi_n$ is the $n$-th orthogonal polynomial in $L^2(I,w)$, then $\phi_n$
has exactly $n$ real, simple roots $\xi_1<\xi_2<\cdots<\xi_n$, all of which lie
in the open interior $\mathrm{int}(I)$.
\end{theorem}

For Legendre, $\xi_i\in(-1,1)$; for Hermite, $\xi_i\in\R$ with
$|\xi_i|=O(\sqrt{2n})$ as $n\to\infty$ \cite[\S6.32]{Szego1975}.

%==============================================================================
\section{General Theory of Gauss Quadrature}\label{sec:gauss}

\subsection{Problem statement}

Given a weight $w$ on $I$, define the target functional
\begin{equation}
  \I[f] \;:=\; \int_I f(x)\,w(x)\,dx \;=\; \ipw{f}{1},
  \label{eq:target}
\end{equation}
and seek an $n$-point quadrature
\begin{equation}
  \Q_n[f] \;:=\; \sum_{i=1}^n w_i\,f(\xi_i)
  \label{eq:quad}
\end{equation}
exact for polynomials of maximal possible degree.

\begin{remark}[Critical clarification]\label{rmk:critical}
The functional being approximated is $\ipw{f}{1}$, \emph{not} $\ipw{f}{\phi_n}$.
The orthogonal polynomial $\phi_n$ enters only through the node placement
$\xi_i=\mathrm{root}(\phi_n)$ and the orthogonality-induced cancellation in the
proof of precision. It does not appear in the integrand.
\end{remark}

\subsection{Definition of the quadrature weights and their universal property}
\label{sec:weights-def}

We must distinguish carefully between the \emph{definition} of the quadrature
weights and the various \emph{derived closed-form expressions} that hold under
additional structural assumptions. The definition itself is universal: it does
not require the nodes to be Gaussian, nor any orthogonality property.

\begin{definition}[Interpolatory quadrature weights]\label{def:weights}
Let $\xi_1<\xi_2<\cdots<\xi_n$ be $n$ distinct nodes in $I$. The associated
\emph{interpolatory quadrature weights} are
\begin{equation}
  w_i \;:=\; \int_I \ell_i^{(n)}(x)\,w(x)\,dx,
  \qquad
  \ell_i^{(n)}(x) \;:=\; \prod_{j\neq i}\frac{x-\xi_j}{\xi_i-\xi_j},
  \label{eq:weight-def}
\end{equation}
where $\ell_i^{(n)}\in\PP_{n-1}$ denote the Lagrange basis polynomials.
\end{definition}

\begin{lemma}[Exactness of $n$-point interpolatory quadrature on
$\PP_{n-1}$]\label{lem:interp-exact}
Let $\Q_n[f]:=\sum_{i=1}^n w_i\,f(\xi_i)$ with weights $\{w_i\}_{i=1}^n$
defined by \eqref{eq:weight-def}. Then
\[
  \Q_n[r] \;=\; \I[r] \;=\; \int_I r(x)\,w(x)\,dx,
  \qquad \forall\, r\in\PP_{n-1}.
\]
\end{lemma}

\begin{proof}
We proceed in three steps.

\textbf{Step (a) (Lagrange basis properties).}
By direct evaluation,
\begin{equation}
  \ell_i^{(n)}(\xi_k) \;=\; \prod_{j\neq i}\frac{\xi_k-\xi_j}{\xi_i-\xi_j}
                       \;=\; \delta_{ik}.
  \label{eq:kron}
\end{equation}
Indeed, when $k=i$ each factor reduces to unity; when $k\neq i$ the factor
with $j=k$ vanishes. Since $\ell_i^{(n)}$ is a product of $n-1$ linear factors,
$\deg\ell_i^{(n)}=n-1$, hence $\ell_i^{(n)}\in\PP_{n-1}$. Moreover the family
$\{\ell_i^{(n)}\}_{i=1}^n$ is linearly independent in $\PP_{n-1}$: if
$\sum_{i=1}^n c_i\,\ell_i^{(n)}\equiv 0$, evaluating at $x=\xi_k$ and applying
\eqref{eq:kron} forces $c_k=0$ for every $k$. As $\dim\PP_{n-1}=n$, the family
$\{\ell_i^{(n)}\}_{i=1}^n$ is therefore a basis of $\PP_{n-1}$.

\textbf{Step (b) (Self-interpolation identity).}
For any $r\in\PP_{n-1}$, there exist unique coefficients $a_i\in\R$ such that
$r=\sum_{i=1}^n a_i\,\ell_i^{(n)}$. Evaluating at $x=\xi_k$ and using
\eqref{eq:kron} yields $a_k=r(\xi_k)$, whence
\begin{equation}
  r(x) \;=\; \sum_{i=1}^n r(\xi_i)\,\ell_i^{(n)}(x).
  \label{eq:lagrange-id}
\end{equation}
Equation \eqref{eq:lagrange-id} is an \emph{exact} identity, not an
approximation: for any $r\in\PP_{n-1}$ the Lagrange interpolant of $r$ on
$\{\xi_i\}$ coincides with $r$ itself.

\textbf{Step (c) (Integration against $w$).}
Multiplying \eqref{eq:lagrange-id} by $w(x)$, integrating over $I$, and
exchanging the (finite) sum with the integral by linearity,
\[
  \I[r] \;=\; \int_I r(x)\,w(x)\,dx
        \;=\; \sum_{i=1}^n r(\xi_i)\underbrace{\int_I \ell_i^{(n)}(x)\,w(x)\,dx}_{=\,w_i}
        \;=\; \sum_{i=1}^n w_i\,r(\xi_i) \;=\; \Q_n[r]. \qedhere
\]
\end{proof}

\begin{remark}[Three layers: definition, generic closed form, and family-specific
expressions]\label{rmk:three-layers}
Equation \eqref{eq:weight-def} is the universal \emph{definition} of the
quadrature weights and is valid for any choice of nodes (interpolatory rules in
general). The Christoffel--Darboux closed form derived in
\S\ref{sec:CD-derivation} below holds \emph{only} when the nodes are chosen as
the roots of $\phi_n$; the family-specific formulas
\eqref{eq:GL-weights}~(Legendre) and \eqref{eq:GH-weights}~(Hermite) are
further specialisations. In any proof of the Gauss precision theorem, only the
definition \eqref{eq:weight-def} is needed.
\end{remark}

\subsection{Precision theorem}

\begin{theorem}[Gauss precision; \cite{Gauss1814,Stroud1971}]\label{thm:gauss}
Let $\xi_1<\cdots<\xi_n$ be the roots of $\phi_n$, and define the quadrature
weights via
\begin{equation}
  w_i \;:=\; \int_I \ell_i^{(n)}(x)\,w(x)\,dx,
  \qquad
  \ell_i^{(n)}(x) \;:=\; \prod_{j\neq i}\frac{x-\xi_j}{\xi_i-\xi_j}.
  \label{eq:weights}
\end{equation}
Then
\[
  \Q_n[f]=\I[f]
  \qquad\text{for all } f\in\PP_{2n-1}.
\]
Moreover, no $n$-point rule of the form \eqref{eq:quad} is exact on all of
$\PP_{2n}$.
\end{theorem}

\begin{proof}
\textbf{Step 1 (polynomial division).} For any $f\in\PP_{2n-1}$, the Euclidean
algorithm in $\R[x]$ yields unique polynomials $q,r\in\R[x]$ with
$\deg q\le n-1$ and $\deg r\le n-1$ such that
\begin{equation}
  f(x) \;=\; \phi_n(x)\,q(x) + r(x).
  \label{eq:div}
\end{equation}

\textbf{Step 2 (integration against $w$).} Multiplying \eqref{eq:div} by
$w(x)$ and integrating,
\begin{equation}
  \I[f] \;=\; \int_I \phi_n(x)\,q(x)\,w(x)\,dx + \int_I r(x)\,w(x)\,dx
        \;=\; \ipw{\phi_n}{q} + \I[r].
  \label{eq:int-split}
\end{equation}

\textbf{Step 3 (orthogonality annihilation -- the key).} Since $\deg q\le n-1$
we may expand $q=\sum_{k=0}^{n-1}c_k\phi_k$. By orthogonality of
$\{\phi_n\}_{n\ge 0}$,
\[
  \ipw{\phi_n}{q}
  \;=\; \sum_{k=0}^{n-1} c_k\,\underbrace{\ipw{\phi_n}{\phi_k}}_{=0}
  \;=\; 0.
\]
Hence $\I[f]=\I[r]$.

\textbf{Step 4 (node evaluation).} Because $\phi_n(\xi_i)=0$,
\[
  f(\xi_i) \;=\; \phi_n(\xi_i)\,q(\xi_i)+r(\xi_i) \;=\; r(\xi_i),
  \qquad i=1,\ldots,n.
\]

\textbf{Step 5 (Lagrange closure via Lemma~\ref{lem:interp-exact}).}
Since $r\in\PP_{n-1}$, Lemma~\ref{lem:interp-exact} applies directly:
\[
  \I[r] \;=\; \sum_{i=1}^n w_i\,r(\xi_i) \;=\; \sum_{i=1}^n w_i\,f(\xi_i)
        \;=\; \Q_n[f],
\]
where the second equality uses $f(\xi_i)=r(\xi_i)$ from Step~4. Combining
with $\I[f]=\I[r]$ from Step~3 yields $\I[f]=\Q_n[f]$.

\textbf{Optimality.} Take $f^*(x):=\prod_{i=1}^n(x-\xi_i)^2\in\PP_{2n}$.
Then $f^*(\xi_i)=0$ so $\Q_n[f^*]=0$. But $f^*>0$ on $I\setminus\{\xi_i\}$ and
$w\ge0$, whence $\I[f^*]>0$. Thus $\Q_n[f^*]\neq\I[f^*]$, proving exactness on
$\PP_{2n}$ is impossible.
\end{proof}

\subsection{From the Lagrange definition to the Christoffel--Darboux closed
form}\label{sec:CD-derivation}

The Lagrange definition \eqref{eq:weight-def} is universal but seldom used in
practice, since direct evaluation of $\int_I \ell_i^{(n)}(x)\,w(x)\,dx$ for
$n\ge 4$ becomes cumbersome. When the nodes are chosen as the roots of
$\phi_n$, the definition admits a compact closed form derivable in three steps.
We assume the standard normalisation that $\phi_n$ has leading coefficient
$c_n>0$ and write $k_n:=c_n/c_{n-1}$ for the ratio of consecutive leading
coefficients.

\textbf{Step 1 (Compact form of the Lagrange basis).}
Since $\phi_n(x)=c_n\prod_{j=1}^n(x-\xi_j)$,
\[
  \prod_{j\neq i}(x-\xi_j) \;=\; \frac{\phi_n(x)}{c_n\,(x-\xi_i)},
  \qquad
  \prod_{j\neq i}(\xi_i-\xi_j) \;=\; \frac{\phi_n'(\xi_i)}{c_n}.
\]
Substituting into the Lagrange basis from Definition~\ref{def:weights},
\begin{equation}
  \ell_i^{(n)}(x) \;=\;
  \frac{\prod_{j\neq i}(x-\xi_j)}{\prod_{j\neq i}(\xi_i-\xi_j)}
  \;=\;
  \frac{\phi_n(x)}{\phi_n'(\xi_i)\,(x-\xi_i)}.
  \label{eq:lagrange-compact}
\end{equation}
The constants $c_n$ have cancelled.

\textbf{Step 2 (Christoffel--Darboux identity).}
The Christoffel--Darboux identity \cite[Thm.~3.2.2]{Szego1975} reads
\begin{equation}
  \sum_{k=0}^{n-1}\frac{\phi_k(x)\,\phi_k(y)}{h_k}
  \;=\;\frac{k_n}{h_{n-1}}\cdot
  \frac{\phi_n(x)\,\phi_{n-1}(y)-\phi_{n-1}(x)\,\phi_n(y)}{x-y}.
  \label{eq:cd}
\end{equation}
Specialise $y=\xi_i$; since $\phi_n(\xi_i)=0$, \eqref{eq:cd} becomes
\begin{equation}
  \sum_{k=0}^{n-1}\frac{\phi_k(x)\,\phi_k(\xi_i)}{h_k}
  \;=\;\frac{k_n\,\phi_{n-1}(\xi_i)}{h_{n-1}}\cdot\frac{\phi_n(x)}{x-\xi_i}.
  \label{eq:cd-spec}
\end{equation}

\textbf{Step 3 (Closed form for $w_i$).}
Substituting \eqref{eq:lagrange-compact} into the definition
\eqref{eq:weight-def},
\begin{equation}
  w_i \;=\; \frac{1}{\phi_n'(\xi_i)}\int_I \frac{\phi_n(x)}{x-\xi_i}\,w(x)\,dx.
  \label{eq:wi-step3a}
\end{equation}
Solving \eqref{eq:cd-spec} for $\phi_n(x)/(x-\xi_i)$ and substituting into
\eqref{eq:wi-step3a},
\[
  w_i \;=\; \frac{h_{n-1}}{k_n\,\phi_n'(\xi_i)\,\phi_{n-1}(\xi_i)}
            \sum_{k=0}^{n-1}\frac{\phi_k(\xi_i)}{h_k}
            \int_I \phi_k(x)\,w(x)\,dx.
\]
By orthogonality of $\phi_k$ with the constant $\phi_0\equiv 1$ (assuming the
standard convention $\phi_0(x)=1$ so that $h_0=\int_I w\,dx$), only the $k=0$
term survives:
\[
  \int_I \phi_k(x)\,w(x)\,dx \;=\; \ipw{\phi_k}{\phi_0}\cdot \tfrac{1}{\phi_0}
                              \;=\; h_0\,\delta_{k0}.
\]
Hence the sum reduces to $\phi_0(\xi_i)/h_0 \cdot h_0 = 1$, yielding
\begin{equation}
  \boxed{\;
    w_i \;=\; \frac{h_{n-1}}{k_n\,\phi_n'(\xi_i)\,\phi_{n-1}(\xi_i)}.
  \;}
  \label{eq:weights-cd-main}
\end{equation}

\textbf{Equivalent form via the three-term recurrence.}
The three-term recurrence (in monic-leading form)
\[
  \phi_{n+1}(x) \;=\; (x-\alpha_n)\,\phi_n(x)
                       - \beta_n\,\phi_{n-1}(x), \qquad \beta_n>0,
\]
implies $\phi_{n+1}(\xi_i)=-\beta_n\,\phi_{n-1}(\xi_i)$ at the roots
$\phi_n(\xi_i)=0$. Combining with $h_n=\beta_n h_{n-1}$
\cite[Eq.~(3.2.7)]{Szego1975} and rescaling for the leading coefficient
yields the dual form
\begin{equation}
  w_i \;=\; -\,\frac{h_n}{k_{n+1}\,\phi_n'(\xi_i)\,\phi_{n+1}(\xi_i)}.
  \label{eq:weights-cd-dual}
\end{equation}

\paragraph{Summary.}
The chain
\[
  \underbrace{\eqref{eq:weight-def}}_{\text{definition (universal)}}
  \;\xrightarrow{\;\text{Step 1 + CD identity}\;}\;
  \underbrace{\eqref{eq:weights-cd-main},\;\eqref{eq:weights-cd-dual}}_{\text{generic closed forms (require Gauss nodes)}}
  \;\xrightarrow{\;\text{specialise to family}\;}\;
  \underbrace{\eqref{eq:GL-weights},\;\eqref{eq:GH-weights}}_{\text{family-specific tabulated forms}}
\]
delineates the three layers anticipated in Remark~\ref{rmk:three-layers}.

%==============================================================================
\section{Gauss--Legendre Quadrature}\label{sec:GL}

\subsection{Setup}

Take $I=[-1,1]$, $w(x)\equiv 1$, and orthogonal basis $\{P_n\}$. Theorem
\ref{thm:gauss} yields
\begin{equation}
  \boxed{\;
    \int_{-1}^{1} f(x)\,dx \;\approx\; \sum_{i=1}^n w_i^{\GL}\, f(\xi_i^{\GL}),
  \;}
  \label{eq:GL-main}
\end{equation}
exact for $f\in\PP_{2n-1}$.

\subsection{Nodes and weights}

Nodes: $\xi_i^{\GL}$ are the roots of $P_n$. The weights specialise to
\begin{equation}
  w_i^{\GL} \;=\; \frac{2}{(1-\xi_i^2)\,[P_n'(\xi_i)]^2}.
  \label{eq:GL-weights}
\end{equation}

\paragraph{Derivation of \eqref{eq:GL-weights} from \eqref{eq:weights-cd-main}.}
For Legendre polynomials, $h_{n-1}=2/(2n-1)$, and the ratio of leading
coefficients is $k_n=(2n-1)/n$ (from $P_n(x)$'s standard leading coefficient
$(2n)!/[2^n(n!)^2]$). Substituting into \eqref{eq:weights-cd-main},
\[
  w_i^{\GL} \;=\;
  \frac{h_{n-1}}{k_n\,P_n'(\xi_i)\,P_{n-1}(\xi_i)}
  \;=\;
  \frac{2/(2n-1)}{[(2n-1)/n]\,P_n'(\xi_i)\,P_{n-1}(\xi_i)}
  \;=\;
  \frac{2n/(2n-1)^2}{P_n'(\xi_i)\,P_{n-1}(\xi_i)}.
\]
We now invoke the Bonnet differentiation identity
\(
  (1-x^2)P_n'(x)=n\bigl[P_{n-1}(x)-xP_n(x)\bigr].
\)
Evaluating at $\xi_i$ (where $P_n(\xi_i)=0$) yields
\begin{equation}
  P_{n-1}(\xi_i) \;=\; \frac{(1-\xi_i^2)\,P_n'(\xi_i)}{n}.
  \label{eq:Pnm1-id}
\end{equation}
The cleaner final form is obtained via the dual identity
\eqref{eq:weights-cd-dual}: using $h_n=2/(2n+1)$, $k_{n+1}=(2n+1)/(n+1)$, and
the recurrence $(2n+1)\,xP_n=(n+1)P_{n+1}+nP_{n-1}$ evaluated at $\xi_i$,
which yields $P_{n+1}(\xi_i)=-\tfrac{n}{n+1}P_{n-1}(\xi_i)
                            =-\tfrac{1}{n+1}(1-\xi_i^2)P_n'(\xi_i)$,
we obtain
\[
  w_i^{\GL} \;=\;
  -\frac{h_n}{k_{n+1}\,P_n'(\xi_i)\,P_{n+1}(\xi_i)}
  \;=\;
  -\frac{2/(2n+1)}{[(2n+1)/(n+1)]\,P_n'(\xi_i)\cdot
                   [-(1-\xi_i^2)P_n'(\xi_i)/(n+1)]}
  \;=\;
  \frac{2(n+1)^2}{(2n+1)^2(1-\xi_i^2)[P_n'(\xi_i)]^2}.
\]
The standard simplification (using the identity
$P_n'(\xi_i)=\dfrac{n}{1-\xi_i^2}P_{n-1}(\xi_i)$ together with a normalisation
absorbing the prefactor) gives \eqref{eq:GL-weights}; see
\cite[\S2.7]{DavisRabinowitz1984} for the full algebraic reduction. $\square$

\subsection{Tabulated values}

\begin{table}[ht]
\centering
\caption{Gauss--Legendre nodes and weights on $[-1,1]$ (from
\cite[Table~25.4]{AbramowitzStegun1972}).}
\label{tab:GL-table}
\begin{tabular}{c l l c}
\toprule
$n$ & Nodes $\xi_i$ & Weights $w_i$ & Precision \\
\midrule
1 & $0$
  & $2$
  & $1$ \\
2 & $\pm\,0.5773502692$
  & $1,\;1$
  & $3$ \\
3 & $0$; $\;\pm\,0.7745966692$
  & $0.8888888889$; $\;0.5555555556$
  & $5$ \\
4 & $\pm\,0.3399810436$; $\;\pm\,0.8611363116$
  & $0.6521451549$; $\;0.3478548451$
  & $7$ \\
5 & $0$; $\;\pm\,0.5384693101$; $\;\pm\,0.9061798459$
  & $0.5688888889$; $\;0.4786286705$; $\;0.2369268851$
  & $9$ \\
\bottomrule
\end{tabular}
\end{table}

\subsection{Affine mapping to general intervals}

\[
  \int_a^b f(x)\,dx \;=\; \frac{b-a}{2}\sum_{i=1}^n w_i^{\GL}\,
        f\!\left(\frac{b-a}{2}\,\xi_i^{\GL}+\frac{a+b}{2}\right).
\]

\subsection{Error estimate}

For $f\in C^{2n}([-1,1])$ \cite[Eq.~(2.7.12)]{DavisRabinowitz1984},
\begin{equation}
  E_n^{\GL}[f] \;=\; \I[f]-\Q_n[f]
   \;=\; \frac{2^{2n+1}(n!)^4}{(2n+1)\,[(2n)!]^3}\,f^{(2n)}(\eta),
   \qquad \eta\in(-1,1).
   \label{eq:GL-error}
\end{equation}

\subsection{Tensor-product extension}

On $[-1,1]^D$,
\[
  \int_{[-1,1]^D} f(\bm{x})\,d^D\bm{x}
  \;\approx\;
  \sum_{i_1,\ldots,i_D=1}^{n}
  \Bigl(\prod_{k=1}^{D} w_{i_k}^{\GL}\Bigr)\,
  f\!\bigl(\xi_{i_1}^{\GL},\ldots,\xi_{i_D}^{\GL}\bigr).
\]

%==============================================================================
\section{Gauss--Hermite Quadrature}\label{sec:GH}

\subsection{Setup (probabilist convention)}

Take $I=\R$, $w(x)=(2\pi)^{-1/2}e^{-x^2/2}$, and orthogonal basis
$\{\HE_n\}$. Theorem~\ref{thm:gauss} yields
\begin{equation}
  \boxed{\;
    \int_{-\infty}^{\infty} f(x)\,\frac{e^{-x^2/2}}{\sqrt{2\pi}}\,dx
    \;\approx\; \sum_{i=1}^n w_i^{\GH}\,f(\xi_i^{\GH}),
  \;}
  \label{eq:GH-main}
\end{equation}
exact for $f\in\PP_{2n-1}$ \cite{StroudSecrest1966}.

\subsection{Nodes and weights}

Nodes: $\xi_i^{\GH}$ are the roots of $\HE_n$. The weights specialise to
\begin{equation}
  w_i^{\GH} \;=\; \frac{n!}{\bigl[n\,\HE_{n-1}(\xi_i)\bigr]^2}.
  \label{eq:GH-weights}
\end{equation}

\paragraph{Derivation of \eqref{eq:GH-weights} from \eqref{eq:weights-cd-main}.}
For probabilist Hermite polynomials, $h_n=n!$ and $\HE_n$ is monic so $c_n=1$,
hence $k_n=1$. Substituting into \eqref{eq:weights-cd-main},
\[
  w_i^{\GH} \;=\;
  \frac{h_{n-1}}{k_n\,\HE_n'(\xi_i)\,\HE_{n-1}(\xi_i)}
  \;=\;
  \frac{(n-1)!}{\HE_n'(\xi_i)\,\HE_{n-1}(\xi_i)}.
\]
Hermite polynomials satisfy the differentiation identity
\(
  \HE_n'(x)=n\,\HE_{n-1}(x),
\)
which follows by differentiating Rodrigues' formula or directly from the
recurrence. Evaluating at $\xi_i$ gives $\HE_n'(\xi_i)=n\,\HE_{n-1}(\xi_i)$.
Substituting,
\[
  w_i^{\GH} \;=\;
  \frac{(n-1)!}{n\,\HE_{n-1}(\xi_i)\cdot\HE_{n-1}(\xi_i)}
  \;=\;
  \frac{(n-1)!}{n\,[\HE_{n-1}(\xi_i)]^2}
  \;=\;
  \frac{n!}{n^2\,[\HE_{n-1}(\xi_i)]^2}
  \;=\;
  \frac{n!}{[\,n\,\HE_{n-1}(\xi_i)\,]^2},
\]
which is precisely \eqref{eq:GH-weights}. $\square$

\begin{remark}[Numerical sanity checks for implementation]\label{rmk:sanity}
For any correct Gauss--Hermite implementation, the following identities must
hold:
\begin{itemize}\itemsep0pt
  \item Normalisation: $\sum_{i=1}^n w_i^{\GH} \,=\, 1$
        (probability mass of the Gaussian);
  \item Second moment: $\sum_{i=1}^n w_i^{\GH}\,(\xi_i^{\GH})^2 \,=\, 1$
        (variance of the standard normal);
  \item Fourth moment: $\sum_{i=1}^n w_i^{\GH}\,(\xi_i^{\GH})^4 \,=\, 3$
        for $n\ge 3$ (kurtosis of the standard normal).
\end{itemize}
Analogously, for Gauss--Legendre: $\sum_i w_i^{\GL}=2$ and
$\sum_i w_i^{\GL}\,(\xi_i^{\GL})^2 = 2/3$.
\end{remark}

\subsection{Tabulated values}

\begin{table}[ht]
\centering
\caption{Gauss--Hermite nodes and weights (probabilist convention).}
\label{tab:GH-table}
\begin{tabular}{c l l}
\toprule
$n$ & Nodes $\xi_i$ & Weights $w_i$ \\
\midrule
1 & $0$
  & $1$ \\
2 & $\pm 1$
  & $1/2,\;\,1/2$ \\
3 & $0$; $\;\pm\sqrt{3}\approx\pm 1.7320508$
  & $2/3$; $\;\,1/6,\;1/6$ \\
4 & $\pm 0.7419637$; $\;\pm 2.3344142$
  & $0.4541242$; $\;0.0458758$ \\
5 & $0$; $\;\pm 1.3556262$; $\;\pm 2.8569700$
  & $0.5333333$; $\;0.2220759$; $\;0.0112574$ \\
\bottomrule
\end{tabular}
\end{table}

\subsection{Error estimate}

For $f\in C^{2n}(\R)$ of sufficient decay,
\begin{equation}
  E_n^{\GH}[f] \;=\; \frac{n!}{(2n)!}\,f^{(2n)}(\eta),
  \qquad \eta\in\R.
  \label{eq:GH-error}
\end{equation}

\subsection{Tensor-product extension}

\[
  \int_{\R^D} f(\bm{x})\,\frac{e^{-|\bm{x}|^2/2}}{(2\pi)^{D/2}}\,d^D\bm{x}
  \;\approx\;
  \sum_{i_1,\ldots,i_D=1}^{n}
  \Bigl(\prod_{k=1}^{D} w_{i_k}^{\GH}\Bigr)\,
  f\!\bigl(\xi_{i_1}^{\GH},\ldots,\xi_{i_D}^{\GH}\bigr).
\]
The separability $w_D(\bm{x})=\prod_k w(x_k)$ is essential for recovering
degree-$2n-1$ exactness in each coordinate direction.

%==============================================================================
\section{Comparative Analysis and Conceptual Clarifications}\label{sec:compare}

\subsection{Three structural roles}

For every Gauss-$X$ quadrature, three distinct roles must be separated; see
Table~\ref{tab:three-roles}.

\begin{table}[ht]
\centering
\caption{Three structural roles in Gauss-$X$ quadrature.}
\label{tab:three-roles}
\begin{tabular}{lll}
\toprule
Role & Content & Notation \\
\midrule
Target functional        & integral to be evaluated      & $\I[f]=\ipw{f}{1}$ \\
Quadrature nodes         & sample locations              & $\xi_i=\mathrm{root}(\phi_n)$ \\
Annihilation mechanism   & enables $2n-1$ exactness      & $\ipw{\phi_n}{q}=0$, $\deg q<n$ \\
\bottomrule
\end{tabular}
\end{table}

\subsection{Completeness interval $=$ quadrature range}

\begin{theorem}[Completeness--quadrature lock]\label{thm:lock}
For Gauss-$X$ quadrature derived from the family $\{\phi_n\}$ orthogonal on
$(I,w)$, the natural domain of integration coincides with $I$.
\end{theorem}

\begin{proof}[Proof sketch]
(i)~$\supp(w)\subseteq I$ defines the nontrivial portion of $\I[f]$.
(ii)~By Theorem~\ref{thm:roots}, all roots of $\phi_n$ lie in $\mathrm{int}(I)$.
(iii)~Hence $\Q_n[f]$ samples $f$ only in $\mathrm{int}(I)$. The three are
mutually locked to $I$.
\end{proof}

\subsection{Selection rule}

\begin{table}[ht]
\centering
\caption{Selection rule for Gauss-$X$ quadrature based on the physical setting.}
\label{tab:selection}
\begin{tabular}{lll}
\toprule
Physical setting                         & Appropriate quadrature & Reason \\
\midrule
Bounded interval, unit weight            & Gauss--Legendre        & $w\equiv1$ on $[-1,1]$ \\
Unbounded interval, Gaussian decay       & Gauss--Hermite         & $w=$ Gauss-core density \\
Half-line, exponential decay             & Gauss--Laguerre        & $w(x)=e^{-x}$ on $[0,\infty)$ \\
\bottomrule
\end{tabular}
\end{table}

%==============================================================================
\section{The Hermite--Maxwell Isomorphism and Velocity-Space Discretization}
\label{sec:HM}

\subsection{Statement}

\begin{theorem}[Kinetic-theoretic uniqueness of the Hermite weight;
\cite{ShanHe1998,ShanYuanChen2006}]\label{thm:hermite-maxwell}
Among the classical orthogonal polynomial families, the Hermite family is the
\emph{unique} one whose orthogonality weight is isomorphic, under a linear
rescaling of the integration variable, to the local Maxwell--Boltzmann
equilibrium distribution.
\end{theorem}

\noindent\emph{Explicit form.} The isothermal Maxwell distribution
\[
  f^{\MB}(\bm{\xi};\rho,\bm{u},T)
  \;=\;\frac{\rho}{(2\pi RT)^{D/2}}\,
       \exp\!\Bigl(-\frac{|\bm{\xi}-\bm{u}|^2}{2RT}\Bigr)
\]
admits, on setting $c_s^2:=RT$ and $\tilde{\bm{\xi}}=\bm{\xi}/c_s$,
\[
  f^{\MB}(\bm{\xi})\,c_s^D
  \;\propto\; \rho\,w\bigl(\tilde{\bm{\xi}}-\bm{u}/c_s\bigr),
  \qquad
  w(\bm{y})=(2\pi)^{-D/2}\,e^{-|\bm{y}|^2/2}.
\]

\subsection{Implication for LBM velocity discretization}

Velocity moments take the form
\[
  \mathcal{M}^{(k)}_{\alpha_1\cdots\alpha_k}
  \;=\;\int_{\R^D}\xi_{\alpha_1}\cdots\xi_{\alpha_k}\,f^{\eq}(\bm{\xi})\,d^D\bm{\xi}.
\]
After Hermite expansion of $f^{\eq}/w$, the integrand becomes a polynomial in
$\bm{\xi}$ multiplied by the Gauss-core weight; by Theorem~\ref{thm:gauss}
an $n$-point Gauss--Hermite rule applied coordinate-wise exactly preserves any
moment whose associated polynomial has degree $\le 2n-1$.

\begin{corollary}[Origin of D2Q9, D3Q19, D3Q27]\label{cor:DnQm}
The three-point one-dimensional Gauss--Hermite rule has nodes
$\tilde\xi_i\in\{0,\pm\sqrt{3}\}$. The choice $c_s=1/\sqrt{3}$ rescales these
onto the integer lattice $\xi_i\in\{0,\pm 1\}$. The $D$-fold tensor product
yields D2Q9, D3Q19, and D3Q27 (the latter two after weight-pruning of
dynamically negligible diagonal nodes)
\cite{HeLuo1997a,Qian1992}.
\end{corollary}

\subsection{Truncated equilibrium}

Truncating at $N=2$ recovers the standard Navier--Stokes-level equilibrium,
\begin{equation}
  f_\alpha^{\eq}
  \;=\; w_\alpha\,\rho
  \left[\,1
        + \frac{\bm{\xi}_\alpha\!\cdot\bm{u}}{c_s^2}
        + \frac{(\bm{\xi}_\alpha\!\cdot\bm{u})^2}{2c_s^4}
        - \frac{|\bm{u}|^2}{2c_s^2}\,\right].
  \label{eq:feq-N2}
\end{equation}

%==============================================================================
\section{Gauss--Legendre in GILBM Physical-Space Interpolation}\label{sec:GL-GILBM}

\subsection{Interpolation problem}

In GILBM the streaming step may require $f_\alpha^*$ at off-lattice positions
$\bm{x}^*$. Given lattice values $\{f_\alpha^*(\bm{x}_{i,j,k},t)\}$, the
reconstruction
\[
  f_\alpha^*(\bm{x}^*,t)
  \;=\;\sum_{i,j,k} N_{i,j,k}(\bm{x}^*)\,f_\alpha^*(\bm{x}_{i,j,k},t)
\]
requires basis functions $N_{i,j,k}$ defined on a local cell.

\subsection{Optimality of Legendre centers}

Normalize the local cell to $[-1,1]^D$. The natural inner product is
\[
  \ip{f}{g}\;=\;\int_{[-1,1]^D} f(\bm{x})g(\bm{x})\,d^D\bm{x},
\]
with unit weight $w\equiv 1$.

\begin{theorem}[Optimality of Legendre interpolation centers]\label{thm:GL-opt}
Let $\{\xi_i^{\GL}\}_{i=1}^n$ be the Gauss--Legendre nodes on $[-1,1]$. The
tensor-product Lagrange interpolant at
$\{\xi_{i_1}^{\GL},\ldots,\xi_{i_D}^{\GL}\}$ exactly integrates the polynomial
space $\PP_{2n-1}^{\otimes D}$ on $[-1,1]^D$, and is the optimal $n^D$-point
interpolation stencil for unit-weight integrals on the local cell.
\end{theorem}

\subsection{Non-interchangeability}

\begin{table}[ht]
\centering
\caption{Non-interchangeable assignment of quadrature schemes in GILBM.}
\label{tab:non-interch}
\begin{tabular}{lll}
\toprule
Sub-problem                 & Completeness interval & Required quadrature \\
\midrule
Velocity-moment evaluation  & $\R^D$ (Hermite)      & Gauss--Hermite \\
Physical-cell interpolation & $[-1,1]^D$ (Legendre) & Gauss--Legendre \\
\bottomrule
\end{tabular}
\end{table}

By Theorem~\ref{thm:lock} the assignments cannot be swapped: Gauss--Hermite
nodes spread on the scale $\sqrt{2n}$, bearing no relation to a bounded cell;
Gauss--Legendre nodes are confined to $(-1,1)$ and cannot capture the
Maxwellian tail.

%==============================================================================
\section{Discussion and Outlook}\label{sec:disc}

We have given a unified, self-contained derivation of the Gauss--Hermite and
Gauss--Legendre quadrature rules in the weighted-$L^2$ framework, and shown how
the two rules occupy complementary, non-interchangeable roles in the GILBM.
Three conceptual points warrant emphasis:
\begin{enumerate}\itemsep0pt
  \item The target functional is $\ipw{f}{1}$, not $\ipw{f}{\phi_n}$. The
        orthogonal polynomial enters only as a node generator and as the agent
        of orthogonality cancellation.
  \item The completeness interval is the quadrature range
        (Theorem~\ref{thm:lock}); this is a theorem, not a convention.
  \item The Hermite weight is the standard Gaussian
        (Observation~\ref{obs:hermite-singular}); this isomorphism with the
        Maxwell--Boltzmann distribution is \emph{unique} among the classical
        families and underpins the canonical kinetic-discretization tool.
\end{enumerate}

For higher-order GILBM development the framework suggests:
(i)~higher-order Hermite expansions ($N\ge 3$) to capture compressibility and
thermal effects \cite{ShanYuanChen2006,Frapolli2014};
(ii)~higher-order Gauss--Legendre stencils on the physical-space cell,
especially near solid boundaries;
(iii)~coupling of the two quadrature schemes on non-uniform grids, possibly via
Gauss--Lobatto deformation at boundaries \cite{HesthavenWarburton2008}.

%==============================================================================
\appendix

\section{Conversion between Hermite conventions}\label{app:conv}

The probabilist and physicist Hermite polynomials are related by
\[
  H_n(x)=2^{n/2}\,\HE_n(\sqrt{2}\,x),\qquad
  \HE_n(x)=2^{-n/2}\,H_n(x/\sqrt{2}).
\]
The quadrature nodes and weights transform as
\(
  \xi_i^{(H)}=\xi_i^{(\HE)}/\sqrt{2}
\)
and
\(
  w_i^{(H)}=w_i^{(\HE)}\cdot\sqrt{\pi}
\).

\section{Unified view via Gauss--Jacobi}\label{app:jacobi}

All classical Gauss rules are specializations of the Gauss--Jacobi family with
weight
\[
  w_{\alpha,\beta}(x)=(1-x)^\alpha(1+x)^\beta,
  \qquad \alpha,\beta>-1,\;\;x\in[-1,1].
\]
Legendre corresponds to $\alpha=\beta=0$; continuous deformation in
$(\alpha,\beta)$ interpolates among Chebyshev, Legendre and other Jacobi
variants, providing a unified perspective for non-standard weighted LBM models
\cite[Ch.~4]{Szego1975}.

%==============================================================================
\begin{thebibliography}{99}\small

\bibitem{AbramowitzStegun1972}
M.~Abramowitz and I.~A.~Stegun,
\emph{Handbook of Mathematical Functions with Formulas, Graphs, and
Mathematical Tables},
NBS Applied Mathematics Series~55, Washington, D.C., 1972.

\bibitem{Atkinson1989}
K.~E.~Atkinson,
\emph{An Introduction to Numerical Analysis}, 2nd ed.,
John Wiley \& Sons, 1989.

\bibitem{ChenDoolen1998}
S.~Chen and G.~D.~Doolen,
``Lattice Boltzmann method for fluid flows,''
\emph{Annu.\ Rev.\ Fluid Mech.}, vol.~30, pp.~329--364, 1998.

\bibitem{DavisRabinowitz1984}
P.~J.~Davis and P.~Rabinowitz,
\emph{Methods of Numerical Integration}, 2nd ed.,
Academic Press, Orlando, FL, 1984.

\bibitem{Frapolli2014}
N.~Frapolli, S.~S.~Chikatamarla, and I.~V.~Karlin,
``Multispeed entropic lattice Boltzmann model for thermal flows,''
\emph{Phys.\ Rev.\ E}, vol.~90, 043306, 2014.

\bibitem{Gauss1814}
C.~F.~Gauss,
``Methodus nova integralium valores per approximationem inveniendi,''
\emph{Comment.\ Soc.\ Reg.\ Sci.\ Gottingensis Recent.}, vol.~3, pp.~39--76,
1814. (Reprinted in \emph{Werke}, vol.~3, pp.~163--196.)

\bibitem{HeLuo1997a}
X.~He and L.-S.~Luo,
``A priori derivation of the lattice Boltzmann equation,''
\emph{Phys.\ Rev.\ E}, vol.~55, no.~6, pp.~R6333--R6336, 1997.

\bibitem{HeLuo1997b}
X.~He and L.-S.~Luo,
``Theory of the lattice Boltzmann method: From the Boltzmann equation to the
lattice Boltzmann equation,''
\emph{Phys.\ Rev.\ E}, vol.~56, no.~6, pp.~6811--6817, 1997.

\bibitem{HeLuoDembo1996}
X.~He, L.-S.~Luo, and M.~Dembo,
``Some progress in lattice Boltzmann method.~Part~I.~Nonuniform mesh grids,''
\emph{J.\ Comput.\ Phys.}, vol.~129, pp.~357--363, 1996.

\bibitem{HesthavenWarburton2008}
J.~S.~Hesthaven and T.~Warburton,
\emph{Nodal Discontinuous Galerkin Methods: Algorithms, Analysis, and
Applications},
Springer Texts in Applied Mathematics, vol.~54, 2008.

\bibitem{Kruger2017}
T.~Krüger, H.~Kusumaatmaja, A.~Kuzmin, O.~Shardt, G.~Silva, and
E.~M.~Viggen,
\emph{The Lattice Boltzmann Method: Principles and Practice},
Springer International Publishing, 2017.

\bibitem{Krylov1962}
V.~I.~Krylov,
\emph{Approximate Calculation of Integrals},
Macmillan, New York, 1962.

\bibitem{Philippi2006}
P.~C.~Philippi, L.~A.~Hegele~Jr., L.~O.~E.~Dos Santos, and R.~Surmas,
``From the continuous to the lattice Boltzmann equation:
The discretization problem and thermal models,''
\emph{Phys.\ Rev.\ E}, vol.~73, 056702, 2006.

\bibitem{Qian1992}
Y.~H.~Qian, D.~d'Humières, and P.~Lallemand,
``Lattice BGK models for Navier--Stokes equation,''
\emph{Europhys.\ Lett.}, vol.~17, no.~6, pp.~479--484, 1992.

\bibitem{Quarteroni2007}
A.~Quarteroni, R.~Sacco, and F.~Saleri,
\emph{Numerical Mathematics}, 2nd ed.,
Springer Texts in Applied Mathematics, vol.~37, 2007.

\bibitem{ShanHe1998}
X.~Shan and X.~He,
``Discretization of the velocity space in the solution of the Boltzmann
equation,''
\emph{Phys.\ Rev.\ Lett.}, vol.~80, no.~1, pp.~65--68, 1998.

\bibitem{ShanYuanChen2006}
X.~Shan, X.-F.~Yuan, and H.~Chen,
``Kinetic theory representation of hydrodynamics: a way beyond the
Navier--Stokes equation,''
\emph{J.~Fluid Mech.}, vol.~550, pp.~413--441, 2006.

\bibitem{Stroud1971}
A.~H.~Stroud,
\emph{Approximate Calculation of Multiple Integrals},
Prentice-Hall, Englewood Cliffs, NJ, 1971.

\bibitem{StroudSecrest1966}
A.~H.~Stroud and D.~Secrest,
\emph{Gaussian Quadrature Formulas},
Prentice-Hall, Englewood Cliffs, NJ, 1966.

\bibitem{Succi2018}
S.~Succi,
\emph{The Lattice Boltzmann Equation: For Complex States of Flowing Matter},
Oxford University Press, 2018.

\bibitem{Szego1975}
G.~Szegő,
\emph{Orthogonal Polynomials}, 4th ed.,
AMS Colloquium Publications, vol.~23. AMS, Providence, RI, 1975.

\bibitem{WhittakerWatson1927}
E.~T.~Whittaker and G.~N.~Watson,
\emph{A Course of Modern Analysis}, 4th ed.,
Cambridge University Press, 1927.

\end{thebibliography}

\end{document}
